\begingroup
\setlength{\LTcapwidth}{\textwidth}
\small
\setlength{\tabcolsep}{4pt}
\captionsetup{font=normalsize}
\begin{longtable}{@{}p{3.0cm}p{5.0cm}p{5.0cm}@{}}
\caption{Posisi penelitian terhadap penelitian terdahulu.}
\label{tab:tab2.6_posisi_penelitian} \\
\toprule
\textbf{Dimensi} & \textbf{Penelitian terdahulu} & \textbf{Penelitian ini} \\
\midrule
\endfirsthead

\multicolumn{3}{@{}l}{\small\emph{Tabel \thetable{} (lanjutan)}} \\
\toprule
\textbf{Dimensi} & \textbf{Penelitian terdahulu} & \textbf{Penelitian ini} \\
\midrule
\endhead

\bottomrule
\endlastfoot

Data &
CGM, multimoda, dokumen klinis, atau data kesehatan terintegrasi &
Logbook dan kanal glukosa yang tersedia pada dataset OhioT1DM \\[5pt]

Keluaran prediksi &
Umumnya berupa nilai atau estimasi kadar glukosa &
Nilai prediksi beserta representasi keadaan yang digunakan untuk conditioning retrieval \\[5pt]

Dasar pembentukan kueri \textit{retrieval} &
Pertanyaan pengguna, kondisi saat ini, atau kebutuhan model bahasa &
Keadaan glukosa yang diprediksi \\[5pt]

Retrieval &
Prediksi glukosa dan retrieval klinis umumnya dibahas sebagai fungsi terpisah; RAG klinis menggunakan pertanyaan atau kondisi saat ini &
Retrieval pedoman dikondisikan oleh keadaan fisiologis masa depan yang diprediksi \\[5pt]

\textit{Grounding} &
Dokumen klinis atau sumber pengetahuan dengan karakteristik yang bervariasi &
Korpus pedoman klinis terkurasi dengan keterlacakan sumber \\[5pt]

Alur keputusan &
Bervariasi, termasuk keluaran model atau bantuan interpretasi &
\textit{Doctor-mediated}; rekomendasi ditinjau oleh tenaga kesehatan \\[5pt]

Infrastruktur &
Sebagian pendekatan memerlukan pemantauan kontinu, integrasi rekam medis, atau representasi pengetahuan formal &
Target lingkungan komputasi ringan dengan representasi pasien berbasis logbook
\\
\end{longtable}
\endgroup
